\chapter{Sensores: Percebendo o Mundo}

\begin{objetivos}
\begin{itemize}
\item Descobrir quando um ator toca em algo ou em uma cor;
\item Ler o teclado e a posição do mouse;
\item Fazer perguntas ao jogador e usar a resposta.
\end{itemize}
\end{objetivos}

\section{Os blocos azul-claros}

Os blocos \bloco[scSensing]{Sensores} deixam o programa \textbf{perceber} o que acontece: se dois atores se tocam, qual tecla está apertada, onde está o mouse. Eles têm formato de \textbf{hexágono} ou \textbf{redondo}, porque \textbf{respondem} a uma pergunta ou entregam um valor.

\begin{pensaassim}[Pense assim: os sentidos]
Você usa os olhos, o tato e os ouvidos para perceber o mundo. Os sensores são os \textbf{sentidos do seu personagem}: ``estou tocando na parede?'' é o tato, ``qual tecla apertaram?'' é a audição. Sem sensores, o personagem não percebe nada ao seu redor.
\end{pensaassim}

\begin{center}
\renewcommand{\arraystretch}{1.8}
\begin{tabular}{@{}ll@{}}
\booleano{tocando em \menu{borda}?} & sim ou não: está tocando? \\
\booleano{tocando na cor \menu{cor}?} & sim ou não: está sobre essa cor? \\
\booleano{tecla \menu{espaço} pressionada?} & sim ou não: a tecla está apertada? \\
\booleano{mouse pressionado?} & sim ou não: o botão está apertado? \\
\variavel{posição x do mouse} & número: onde o mouse está \\
\end{tabular}
\end{center}

\begin{vocabulario}[Vocabulário: verdadeiro e falso]
Os blocos hexagonais respondem sempre \textbf{verdadeiro} (sim) ou \textbf{falso} (não). Essas respostas se chamam \textbf{valores booleanos}.
\end{vocabulario}

\section{Usando um sensor com o bloco se}

Um sensor sozinho não faz nada: ele precisa de um bloco que \textbf{decida}. O mais comum é o \bloco[scControl]{se \booleano{...} então}. Vamos aprender a fundo no capítulo de Condições, mas já podemos usar:

\begin{codigo}[Gato assustado]
\begin{pilha}
\bk[scEvents]{quando \bandeira\ for clicado}
\sbc[scControl]{sempre}
\sbc[scControl]{se \booleano{tocando em \menu{mouse}?} então}
\bk[scLooks]{diga \entrada{Ai!} por \entrada{0.5} segundos}
\sbcend
\sbcend
\end{pilha}
\end{codigo}

\begin{dica}
Veja que o \bloco[scControl]{se} fica \textbf{dentro} do \bloco[scControl]{sempre}. O ``sempre'' vigia o tempo todo, e o ``se'' decide o que fazer a cada olhada.
\end{dica}

\section{Perguntando ao jogador}

Com \bloco[scSensing]{pergunte \entrada{Qual é o seu nome?} e espere}, o Scratch mostra uma caixinha de texto na tela. O que o jogador digitar fica guardado no bloco \variavel{resposta}.

\begin{missao}[Missão 8: O gatinho curioso]
Faça o gatinho perguntar o nome do jogador e depois cumprimentar usando a resposta.
\end{missao}

\begin{codigo}[Perguntar e cumprimentar]
\begin{pilha}
\bk[scEvents]{quando \bandeira\ for clicado}
\bk[scSensing]{pergunte \entrada{Qual é o seu nome?} e espere}
\bk[scLooks]{diga \reporter{junte \entrada{Olá, } \variavel{resposta}} por \entrada{3} segundos}
\end{pilha}
\end{codigo}

O bloco \bloco[scOperators]{junte \entrada{Olá, } \variavel{resposta}} (da categoria Operadores) cola dois textos em um só. Se o jogador digitar ``Ana'', o gatinho diz ``Olá, Ana''.

\begin{cuidado}
Não esqueça do espaço depois da vírgula em \entrada{Olá, }, senão o gatinho vai dizer ``Olá,Ana'' tudo junto.
\end{cuidado}

\section{Hora de praticar}

\begin{exercicio}[Qual é a sua idade?]
Faça o gatinho perguntar a idade do jogador e depois dizer a resposta por 2 segundos.
\end{exercicio}

\begin{solucao}
\begin{pilha}
\bk[scEvents]{quando \bandeira\ for clicado}
\bk[scSensing]{pergunte \entrada{Qual é a sua idade?} e espere}
\bk[scLooks]{diga \variavel{resposta} por \entrada{2} segundos}
\end{pilha}
\end{solucao}

\begin{exercicio}[Siga o mouse]
Faça o gatinho \textbf{ir para o ponteiro do mouse} o tempo todo.
\end{exercicio}

\begin{solucao}
\begin{pilha}
\bk[scEvents]{quando \bandeira\ for clicado}
\sbc[scControl]{sempre}
\bk[scMotion]{vá para \menu{ponteiro do mouse}}
\sbcend
\end{pilha}
\end{solucao}

\begin{exercicio}[Bateu na parede]
Faça o gatinho andar sempre, e dizer ``Ui!'' sempre que tocar na borda do palco.
\end{exercicio}

\begin{solucao}
\begin{pilha}
\bk[scEvents]{quando \bandeira\ for clicado}
\sbc[scControl]{sempre}
\bk[scMotion]{mova \entrada{5} passos}
\sbc[scControl]{se \booleano{tocando em \menu{borda}?} então}
\bk[scLooks]{diga \entrada{Ui!} por \entrada{0.5} segundos}
\bk[scMotion]{gire \entrada{180} graus}
\sbcend
\sbcend
\end{pilha}
\end{solucao}

\begin{desafio}[Adivinha a cor]
Faça o gatinho dizer ``Estou na cor vermelha!'' sempre que estiver tocando em uma área vermelha do cenário. Dica: use \booleano{tocando na cor \menu{cor}?} e escolha a cor com o conta-gotas.
\end{desafio}

\begin{solucao}
\begin{pilha}
\bk[scEvents]{quando \bandeira\ for clicado}
\sbc[scControl]{sempre}
\sbc[scControl]{se \booleano{tocando na cor \menu{vermelho}?} então}
\bk[scLooks]{diga \entrada{Estou na cor vermelha!}}
\sbcend
\sbcend
\end{pilha}
Depois de mover o gatinho para longe do vermelho, o balão continua na tela. Para limpá-lo, acrescente um \bloco[scLooks]{diga \entrada{}} quando ele sair.
\end{solucao}
